\section{Discussion}
\label{sec:discussion}

\emph{Interpretation.} We read the results through two claims. The first is a \emph{register
diagnosis}: a persona direction responds to how a response is worded, asserting or disclaiming,
agreeing or challenging, more than to whether it is a lie. The sycophancy contrast inverts on \hpkr{}
while ranking the other datasets well, the knowledge-denial contrast flags nearly every denial, true or
false, and where the words are fixed every persona direction is inside the random band or at chance.
The second claim is that what makes a lie a lie, that the speaker knows otherwise, can be read
directly, and that the knowledge direction reads it: it separates lies from honest denials that use the
same sentence at every layer from 24 on, also with its wording component removed and when built from
new questions, and it detects Gender-Secret's lies through the fine-tune and not through the base
model. Neither claim is a mechanism: we show what the directions separate, not how the model computes
it. The steering result fits the register reading. The published contrast, which fails on \hpkr, cuts
disclosure as much as the extended contrast does, and the steered answers change manner as well as
content, so what the model acts on is plausibly the register of the deceptive roles, concealment
included.

\emph{Model organisms.} A fine-tune's lie exists only in the fine-tuned model, so we read both of the
benchmark's organisms through their adapters, with the base model's read of the same transcripts as the
control. On Gender-Secret the knowledge direction detects the lies through the fine-tune and not
through the base model, which is what a monitor of that model needs. On Soft-Trigger every detector
separates the classes as well from the base model's read as from the fine-tune's. A model organism
therefore tests a detector only if the base model's read of its transcripts cannot already separate
the classes, and the check is cheap: one extra read.

\emph{Limitations.}
\begin{itemize}
\item \emph{Strict thresholds.} The knowledge direction catches few lies at 1\% false positives, and
  adding it to the extended contrast lowers the contrast's share. The selected detector is a ranker
  first.
\item \emph{A supervised probe keeps up on one slice.} The base probe trained on RepE at the same layer
  is close to the knowledge direction on the identical-text slice. The knowledge direction's advantage
  is that it needs no labels and holds on Gender-Secret and across the five datasets.
\item \emph{Designed after the failure.} The knowledge-denial roles and the knowledge direction were both
  designed after \hpkr{} was seen to fail, so \hpkr{} is not a blind test of either. Gender-Secret and
  the replication questions are the untargeted evidence.
\item \emph{Random bands and one fine-tuned test.} On several datasets random directions reach AUROCs of
  0.9 or more, and on Gender-Secret above 0.80, so single-dataset margins are small. Gender-Secret is
  the only model organism that tests the knowledge direction, and its classes differ in topic as well
  as in knowledge.
\item \emph{What the knowledge direction reads.} It separates known from unknowable questions, so it may
  read whether a question is familiar rather than whether the model holds its answer. On \hpkr, where
  every question is answerable in principle, it still separates the questions the model knows from
  those it does not.
\item \emph{Statistics and scope.} Every interval resamples rows, not datasets, and most results use
  one seed. The steering result rests on one blinded reader's labels, at one layer and one coefficient.
  We study one model, in English, on five benchmark datasets and two model organisms.
\end{itemize}
